\documentclass[11pt]{article}
\usepackage[]{bm}
% basic packages
\usepackage[margin=1in]{geometry}
\usepackage[pdftex]{graphicx}
\usepackage{amsmath,amssymb,amsthm}
\usepackage{custom}
\usepackage{lipsum}

% page formatting
\usepackage{fancyhdr}
\pagestyle{fancy}

\renewcommand{\sectionmark}[1]{\markright{\textsf{\arabic{section}. #1}}}
\renewcommand{\subsectionmark}[1]{}
\lhead{\textbf{\thepage} \ \ \nouppercase{\rightmark}}
\chead{}
\rhead{}
\lfoot{}
\cfoot{}
\rfoot{}
\setlength{\headheight}{14pt}

\linespread{1.03} % give a little extra room
\setlength{\parindent}{0.2in} % reduce paragraph indent a bit
\setcounter{secnumdepth}{2} % no numbered subsubsections
\setcounter{tocdepth}{2} % no subsubsections in ToC
\newcommand{\hb}{\hbar}
\renewcommand{\vec}{\mathbf}
\begin{document}

% make title page
\thispagestyle{empty}
\bigskip \
\vspace{0.1cm}

\begin{center}
{\fontsize{22}{22} \selectfont Solid State Physics}
\vskip 16pt
{\fontsize{36}{36} \selectfont \bf \sffamily PHYS 412}
\vskip 24pt
{\fontsize{18}{18} \selectfont \rmfamily Sabit} 
\vskip 6pt
{\fontsize{14}{14} \selectfont \ttfamily ass15@rice.edu} 
\vskip 24pt
\end{center}

{\parindent0pt \baselineskip=15.5pt 
%write something here 



% make table of contents
%\newpage
%\microtoc
\newpage

% main content
\section*{Electrons in Periodic Potential}
12 March, 2025
\begin{itemize}
	\item Previously we have looked at Fermi Gas, Sommerfeld Model (Drude picture in Fermi Gas) but we didn't consider any lattice system. Tight Binding showed Bonds $\to $ energy bands, and there we get the notion of ``periodicity" with particular band gap features showing at the spatial frequencies $\pi / a$. When we consider lattice, the Reciprocal Lattice is an important idea. For instance Bragg Reflection of waves happens when change in a wavevector is a reciprocal lattice vector. 

	\item The lattice problem can be treated as a perturbation to the general problem. Solution to unperturbed and perturbed system (where $
			H = H_0 + \lambda V_1$, unperturbed and perturbed are $H_0$ and $\lambda V_1$ respectively).
		\begin{align*}
			H_0 \psi_0(x) ^{(n)} &= E_0 ^{(n)} \psi_0(x)^{(n)}\\
			H \psi(x) &= E(x)  
		\end{align*}
We have solved the problem of finding the energy corrections (or perturbations) in Perturbation Theory for Quantum Mechanics (PHYS 312). 
\[
E_1 = \bra{\psi_0} V \ket{\psi_0}
\] 
\[
	\ket{\psi_1} = \sum_{m \neq  n} \frac{\bra{\psi^{m} } V \ket{\psi^{n}} }{E^{(n)} - E^{(m)} } \ket{\psi^{
	(m)}}
\] 

\item Fermi gas is the unperturbed starting point. Here the Hamiltonian looks like the Free Particle.
	\[
	H_0 = \frac{p^2}{2m} \quad \varepsilon_0(k) = \frac{\hb ^2 k^2 }{2m}
	\] 
	The states are $\ket{\vec{k}} = \frac{1}{\sqrt{L^3} } e^{i \vec{k} \cdot  \vec{r}}$. Now invoke a potential $V(r)$ that has the property $V(\vec{r}+ \vec{R})$. The matrix element of the perturbation is given by 
	\[
	\bra{\vec{k} '} V \ket{\vec{k}} = \int \mathrm{d} \vec{r} \, 
\frac{e^{ - i \vec{k}' \cdot  \vec{r}}}{\sqrt{L^3}} V(\vec{r})
\frac{e^{  i \vec{k} \cdot  \vec{r}}}{\sqrt{L^3}} = \frac{1}{L^3} \int \mathrm{d} \vec{r} \, 
	e^{ - i (\vec{k} ' - \vec{k} ) \cdot  \vec{r}} V(\vec{r}) = V_{\vec{k}' - \vec{k}}
	\] 
This given integral is zero unless $\vec{k}  - \vec{k} '  = \vec{G}$. 

\item Any function with real-space periodicity of the lattice can be built out of an appropriate combination of reciprocal lattice vectors
	\[
		f(\vec{r}) = f(\vec{r} + \vec{R}) = \sum_{\vec{G}} f_{\vec{G}} e^{i \vec{G} \cdot  \vec{r}} 
	\] 
	The non-vanishing Fourier components of the lattice potentials are the reciprocal lattice vectors.  
\item General Perturbation Theory result 
	\[
	\epsilon(\vec{k}) = 
	e_0(\vec{k}) + 
	\bra{\vec{k}} V \ket{\vec{k}} + 
	\sum_{\vec{k} \neq \vec{k}'} 
	\frac{
| \bra{\vec{k}'} V \ket{\vec{k}} | ^2
	}{\epsilon_0(\vec{k}) - \epsilon_0 (\vec{k}')}
	\]
	Now we want to take the sum like $\vec{k} ' = \vec{k} + \vec{G}$ for non-zero $\vec{G}$ but this blows up when we reach near the edge of the lattice ($\epsilon_0(\vec{k}) = \epsilon_0(\vec{k}')$). Which is a motivation to use Degenerate Perturbation Theory at the zone boundary. 
\item We define the matrix elements $\bra{\vec{k}'} V \ket{\vec{k}'} , \bra{\vec{k}} V \ket{\vec{k}} , \bra{\vec{k}'} V \ket{\vec{k}} , \bra{\vec{k}} V \ket{\vec{k}'}$ such that
	\begin{align*}
	\langle \vec{k} | H | \vec{k} \rangle &= \epsilon_0(\vec{k}) \\
	\langle \vec{k}'  | H | \vec{k}' \rangle &= \epsilon_0(\vec{k}') = \epsilon_0(\vec{k}+\vec{G}) \\
	\langle \vec{k} | H|\vec{k}' \rangle &= V_{\vec{k} - \vec{k}'} = V_G^{*} \\
	\langle \vec{k}' | H| \vec{k} \rangle &= V_G 
	\end{align*}

\item The bigger the Fourier Gap the bigger the energy band gap (?) 
\item In one 1D $V(x) = V_0 \cos (2 \pi x / a)$, at zone boundary we will get 
\item Perturbed wavefunctions always look like Bloch Waves. 

\item Blochs Theorem
\item Bloch States are not momentum eigenstates. $\hb \vec{k}$ is the crystal momentum. This can be observed by operating the usual momentum operator $\vec{\mathcal P} = - i \hb \nabla $
\end{itemize}




















\end{document}
